\documentclass[12pt]{exam}
\usepackage[utf8]{inputenc}
\usepackage{amsmath}
\usepackage{newpxtext,newpxmath}
%\usepackage{libertinus}
\usepackage[T1]{fontenc}
\usepackage{microtype} % Required for proper font rendering
\usepackage{array}

% Formatting for the 'parts' environment
\renewcommand{\thepartno}{\alph{partno}}
\renewcommand{\partlabel}{(\thepartno)}

\begin{document}

% --- VERSION A ---
\noindent \textbf{ECON 1001 — Spring 2026 — Midterm 2 Grading Materials} \\

\noindent Here all versions of the free response problems, for grading reference.

\section*{Version A}
\begin{questions}
    \setcounter{question}{21} % Starts the next question at 22

    \question{Bob and Linda run a diner. They spend their time making Burgers and baking Cookies. The table below shows their maximum production of each good per 8-hour shift.}

    \begin{center}
        \renewcommand{\arraystretch}{1.5}
        \begin{tabular}{|l|c|c|}
            \hline
            \textbf{Worker} & \textbf{Burgers per shift} & \textbf{Cookies per shift} \\
            \hline
            \textbf{Bob}    & 20                         & 60                         \\
            \hline
            \textbf{Linda}  & 40                         & 80                         \\
            \hline
        \end{tabular}
    \end{center}

    \begin{parts}
        \part Who has the absolute advantage in making Burgers? Briefly explain why. Box your answer (for this and each subsequent question).
        \part Compute Bob's opportunity cost of making 1 Burger. Show your work.
        \part Who has the comparative advantage in baking Cookies? Briefly explain why.
        \part Suppose Bob and Linda decide to specialize and trade with one another. What is the range of prices (in terms of cookies) they would agree upon to trade 1 Burger?
    \end{parts}

    \question{The table below shows the Marginal Cost (MC) of producing widgets for three different firms. Widgets are sold in a competitive market for \textbf{\$60 each}. Initially, the government gives each firm exactly \textbf{2 permits}. One permit is required to produce one widget. (Assume firms cannot produce partial widgets.)}

    \begin{center}
        \renewcommand{\arraystretch}{1.5}
        \begin{tabular}{|l|c|c|c|}
            \hline
            \textbf{Widget Number} & \textbf{Firm A's MC} & \textbf{Firm B's MC} & \textbf{Firm C's MC} \\
            \hline
            1st Widget             & \$10                 & \$15                 & \$20                 \\
            \hline
            2nd Widget             & \$25                 & \$30                 & \$50                 \\
            \hline
            3rd Widget             & \$40                 & \$45                 & \$70                 \\
            \hline
        \end{tabular}
    \end{center}

    \begin{parts}
        \part Initially, each firm uses its 2 permits to produce 2 widgets. What is the combined total cost of production for these 6 widgets?
        \part Suppose no permits have been traded yet. What is Firm A's \textbf{maximum} willingness-to-pay (WTP) for a 3rd permit, and what is Firm C's \textbf{minimum} willingness-to-accept (WTA) to sell one of its 2 permits?
        \part Assume transaction costs are zero and firms are allowed to trade permits. After all mutually beneficial trades have occurred, what is the new total cost of production for the 6 widgets?
    \end{parts}
\end{questions}

\newpage

% --- VERSION B ---
\section*{Version B}
\begin{questions}
    \setcounter{question}{21}

    \question{Bob and Linda run a diner. They spend their time making Burgers and baking Cookies. The table below shows their maximum production of each good per 8-hour shift.}

    \begin{center}
        \renewcommand{\arraystretch}{1.5}
        \begin{tabular}{|l|c|c|}
            \hline
            \textbf{Worker} & \textbf{Burgers per shift} & \textbf{Cookies per shift} \\
            \hline
            \textbf{Bob}    & 10                         & 40                         \\
            \hline
            \textbf{Linda}  & 30                         & 90                         \\
            \hline
        \end{tabular}
    \end{center}

    \begin{parts}
        \part Who has the absolute advantage in making Burgers? Briefly explain why. Box your answer (for this and each subsequent question).
        \part Compute Bob's opportunity cost of making 1 Burger. Show your work.
        \part Who has the comparative advantage in baking Cookies? Briefly explain why.
        \part Suppose Bob and Linda decide to specialize and trade with one another. What is the range of prices (in terms of cookies) they would agree upon to trade 1 Burger?
    \end{parts}

    \question{The table below shows the Marginal Cost (MC) of producing widgets for three different firms. Widgets are sold in a competitive market for \textbf{\$60 each}. Initially, the government gives each firm exactly \textbf{2 permits}. One permit is required to produce one widget. (Assume firms cannot produce partial widgets.)}

    \begin{center}
        \renewcommand{\arraystretch}{1.5}
        \begin{tabular}{|l|c|c|c|}
            \hline
            \textbf{Widget Number} & \textbf{Firm A's MC} & \textbf{Firm B's MC} & \textbf{Firm C's MC} \\
            \hline
            1st Widget             & \$12                 & \$16                 & \$22                 \\
            \hline
            2nd Widget             & \$20                 & \$32                 & \$48                 \\
            \hline
            3rd Widget             & \$38                 & \$44                 & \$75                 \\
            \hline
        \end{tabular}
    \end{center}

    \begin{parts}
        \part Initially, each firm uses its 2 permits to produce 2 widgets. What is the combined total cost of production for these 6 widgets?
        \part Suppose no permits have been traded yet. What is Firm A's \textbf{maximum} willingness-to-pay (WTP) for a 3rd permit, and what is Firm C's \textbf{minimum} willingness-to-accept (WTA) to sell one of its 2 permits?
        \part Assume transaction costs are zero and firms are allowed to trade permits. After all mutually beneficial trades have occurred, what is the new total cost of production for the 6 widgets?
    \end{parts}
\end{questions}

\newpage

% --- VERSION C ---
\section*{Version C}
\begin{questions}
    \setcounter{question}{21}

    \question{Bob and Linda run a diner. They spend their time making Burgers and baking Cookies. The table below shows their maximum production of each good per 8-hour shift.}

    \begin{center}
        \renewcommand{\arraystretch}{1.5}
        \begin{tabular}{|l|c|c|}
            \hline
            \textbf{Worker} & \textbf{Burgers per shift} & \textbf{Cookies per shift} \\
            \hline
            \textbf{Bob}    & 25                         & 50                         \\
            \hline
            \textbf{Linda}  & 50                         & 150                        \\
            \hline
        \end{tabular}
    \end{center}

    \begin{parts}
        \part Who has the absolute advantage in making Burgers? Briefly explain why. Box your answer (for this and each subsequent question).
        \part Compute Bob's opportunity cost of making 1 Burger. Show your work.
        \part Who has the comparative advantage in baking Cookies? Briefly explain why.
        \part Suppose Bob and Linda decide to specialize and trade with one another. What is the range of prices (in terms of cookies) they would agree upon to trade 1 Burger?
    \end{parts}

    \question{The table below shows the Marginal Cost (MC) of producing widgets for three different firms. Widgets are sold in a competitive market for \textbf{\$70 each}. Initially, the government gives each firm exactly \textbf{2 permits}. One permit is required to produce one widget. (Assume firms cannot produce partial widgets.)}

    \begin{center}
        \renewcommand{\arraystretch}{1.5}
        \begin{tabular}{|l|c|c|c|}
            \hline
            \textbf{Widget Number} & \textbf{Firm A's MC} & \textbf{Firm B's MC} & \textbf{Firm C's MC} \\
            \hline
            1st Widget             & \$15                 & \$20                 & \$25                 \\
            \hline
            2nd Widget             & \$30                 & \$35                 & \$60                 \\
            \hline
            3rd Widget             & \$45                 & \$50                 & \$80                 \\
            \hline
        \end{tabular}
    \end{center}

    \begin{parts}
        \part Initially, each firm uses its 2 permits to produce 2 widgets. What is the combined total cost of production for these 6 widgets?
        \part Suppose no permits have been traded yet. What is Firm A's \textbf{maximum} willingness-to-pay (WTP) for a 3rd permit, and what is Firm C's \textbf{minimum} willingness-to-accept (WTA) to sell one of its 2 permits?
        \part Assume transaction costs are zero and firms are allowed to trade permits. After all mutually beneficial trades have occurred, what is the new total cost of production for the 6 widgets?
    \end{parts}
\end{questions}

\newpage

% --- VERSION D ---
\section*{Version D}
\begin{questions}
    \setcounter{question}{21}

    \question{Bob and Linda run a diner. They spend their time making Burgers and baking Cookies. The table below shows their maximum production of each good per 8-hour shift.}

    \begin{center}
        \renewcommand{\arraystretch}{1.5}
        \begin{tabular}{|l|c|c|}
            \hline
            \textbf{Worker} & \textbf{Burgers per shift} & \textbf{Cookies per shift} \\
            \hline
            \textbf{Bob}    & 15                         & 60                         \\
            \hline
            \textbf{Linda}  & 20                         & 40                         \\
            \hline
        \end{tabular}
    \end{center}

    \begin{parts}
        \part Who has the absolute advantage in making Burgers? Briefly explain why. Box your answer (for this and each subsequent question).
        \part Compute Bob's opportunity cost of making 1 Burger. Show your work.
        \part Who has the comparative advantage in baking Cookies? Briefly explain why.
        \part Suppose Bob and Linda decide to specialize and trade with one another. What is the range of prices (in terms of cookies) they would agree upon to trade 1 Burger?
    \end{parts}

    \question{The table below shows the Marginal Cost (MC) of producing widgets for three different firms. Widgets are sold in a competitive market for \textbf{\$80 each}. Initially, the government gives each firm exactly \textbf{2 permits}. One permit is required to produce one widget. (Assume firms cannot produce partial widgets.)}

    \begin{center}
        \renewcommand{\arraystretch}{1.5}
        \begin{tabular}{|l|c|c|c|}
            \hline
            \textbf{Widget Number} & \textbf{Firm A's MC} & \textbf{Firm B's MC} & \textbf{Firm C's MC} \\
            \hline
            1st Widget             & \$20                 & \$25                 & \$30                 \\
            \hline
            2nd Widget             & \$30                 & \$35                 & \$65                 \\
            \hline
            3rd Widget             & \$40                 & \$55                 & \$85                 \\
            \hline
        \end{tabular}
    \end{center}

    \begin{parts}
        \part Initially, each firm uses its 2 permits to produce 2 widgets. What is the combined total cost of production for these 6 widgets?
        \part Suppose no permits have been traded yet. What is Firm A's \textbf{maximum} willingness-to-pay (WTP) for a 3rd permit, and what is Firm C's \textbf{minimum} willingness-to-accept (WTA) to sell one of its 2 permits?
        \part Assume transaction costs are zero and firms are allowed to trade permits. After all mutually beneficial trades have occurred, what is the new total cost of production for the 6 widgets?
    \end{parts}
\end{questions}

\newpage

% --- VERSION E ---
\section*{Version E}
\begin{questions}
    \setcounter{question}{21}

    \question{Bob and Linda run a diner. They spend their time making Burgers and baking Cookies. The table below shows their maximum production of each good per 8-hour shift.}

    \begin{center}
        \renewcommand{\arraystretch}{1.5}
        \begin{tabular}{|l|c|c|}
            \hline
            \textbf{Worker} & \textbf{Burgers per shift} & \textbf{Cookies per shift} \\
            \hline
            \textbf{Bob}    & 30                         & 90                         \\
            \hline
            \textbf{Linda}  & 60                         & 120                        \\
            \hline
        \end{tabular}
    \end{center}

    \begin{parts}
        \part Who has the absolute advantage in making Burgers? Briefly explain why. Box your answer (for this and each subsequent question).
        \part Compute Bob's opportunity cost of making 1 Burger. Show your work.
        \part Who has the comparative advantage in baking Cookies? Briefly explain why.
        \part Suppose Bob and Linda decide to specialize and trade with one another. What is the range of prices (in terms of cookies) they would agree upon to trade 1 Burger?
    \end{parts}

    \question{The table below shows the Marginal Cost (MC) of producing widgets for three different firms. Widgets are sold in a competitive market for \textbf{\$50 each}. Initially, the government gives each firm exactly \textbf{2 permits}. One permit is required to produce one widget. (Assume firms cannot produce partial widgets.)}

    \begin{center}
        \renewcommand{\arraystretch}{1.5}
        \begin{tabular}{|l|c|c|c|}
            \hline
            \textbf{Widget Number} & \textbf{Firm A's MC} & \textbf{Firm B's MC} & \textbf{Firm C's MC} \\
            \hline
            1st Widget             & \$5                  & \$10                 & \$15                 \\
            \hline
            2nd Widget             & \$15                 & \$20                 & \$40                 \\
            \hline
            3rd Widget             & \$25                 & \$35                 & \$60                 \\
            \hline
        \end{tabular}
    \end{center}

    \begin{parts}
        \part Initially, each firm uses its 2 permits to produce 2 widgets. What is the combined total cost of production for these 6 widgets?
        \part Suppose no permits have been traded yet. What is Firm A's \textbf{maximum} willingness-to-pay (WTP) for a 3rd permit, and what is Firm C's \textbf{minimum} willingness-to-accept (WTA) to sell one of its 2 permits?
        \part Assume transaction costs are zero and firms are allowed to trade permits. After all mutually beneficial trades have occurred, what is the new total cost of production for the 6 widgets?
    \end{parts}
\end{questions}

\end{document}